\documentclass[11pt]{article}

% ---------------------------------------------------------------------------
% Packages
% ---------------------------------------------------------------------------
\usepackage[margin=1in]{geometry}
\usepackage{amsmath, amssymb}
\usepackage{bbm}
\usepackage{microtype}
\usepackage[hidelinks]{hyperref}

% ---------------------------------------------------------------------------
% Math shortcuts
% ---------------------------------------------------------------------------
\DeclareMathOperator{\Var}{Var}
\DeclareMathOperator*{\argmax}{arg\,max}
\newcommand{\SE}{\operatorname{SE}}
\newcommand{\ind}{\mathbbm{1}}

% A lightweight, non-formal way to set off the worked examples.
\newcommand{\example}[1]{\smallskip\noindent\textit{Example.}\ #1\smallskip}

\title{\textbf{The Probability Behind \emph{Probability Poker}}}
\author{[Your Name] \\ \small First-Year Computer Science, Stanford University}
\date{\today}

\begin{document}

\maketitle

\begin{abstract}
\noindent
\emph{Probability Poker} is a heads-up Texas Hold'em game I built so that the
opponent bot makes every decision from probability theory rather than scripted
poker rules. The bot reasons under uncertainty in four connected steps: it
estimates its chance of winning by Monte Carlo simulation, maintains a Bayesian
belief about my hand, \emph{learns} my personal tendencies from past hands, and
finally chooses the action with the highest expected value. This paper explains
the mathematics behind each step and how they chain together. I write
$\Pr(\cdot)$ for probability and $\mathbb{E}[\cdot]$ for expectation.
\end{abstract}

% ===========================================================================
\section{Monte Carlo Simulation: estimating equity}
% ===========================================================================

The number the bot most wants is its \emph{equity}, $p=\Pr(\text{it wins the
hand})$. Computing $p$ exactly means averaging over every way the hidden cards
could fall, and that space is huge: before the flop the opponent holds one of
$\binom{50}{2}=1225$ possible two-card hands, and for each the remaining board
can come $\binom{48}{5}=1{,}712{,}304$ ways---about $2\times10^{9}$ joint
outcomes. Enumerating that on every decision is hopeless, so I estimate $p$ by
sampling instead.

The justification is the law of large numbers. I run $N$ independent simulated
showdowns; in trial $s$ I deal random cards into the unknown slots (the
opponent's hole cards and the undealt board), score both five-card hands, and
record the indicator
\begin{equation}
X_s=\ind[\text{the bot's best hand beats the opponent's}].
\end{equation}
Each $X_s$ is Bernoulli$(p)$ and they are i.i.d., so the sample mean estimates
$p$ (and the same tally gives losses and ties):
\begin{equation}
\hat p=\frac1N\sum_{s=1}^N X_s=\frac{\text{Wins}}{N},\qquad
\hat q=\frac{\text{Losses}}{N},\qquad
\widehat{\Pr}(\text{tie})=\frac{\text{Ties}}{N}.
\end{equation}
Because $\mathbb{E}[X_s]=p$, the estimator is unbiased: $\mathbb{E}[\hat p]=p$.
By independence its variance is $\Var(\hat p)=p(1-p)/N$, so the standard error is
\begin{equation}
\SE(\hat p)=\sqrt{\frac{p(1-p)}{N}}\;\le\;\frac{1}{2\sqrt N},
\end{equation}
using $p(1-p)\le\tfrac14$. The intuition I find most useful is that the error
falls like $1/\sqrt N$ \emph{no matter how complicated} the outcome space is---
this is exactly why sampling beats enumeration here. To match accuracy to cost I
run more trials when more is unknown (about $7000$ preflop) and fewer when
little is left to chance (about $3000$ on the river). At $N=5000$ the error bound
is roughly $\pm1.4$ percentage points, which is plenty to drive good decisions.

One important detail connects this section to the next: I do \emph{not} draw the
opponent's hand uniformly. I draw it from my current belief about their strength,
so the simulation is weighted toward the hands I think they actually hold.

% ===========================================================================
\section{Bayesian Updating: reading the opponent's hand}\label{sec:bayes}
% ===========================================================================

Since the bot can't see my cards, it tracks a probability distribution over how
strong my hand is, grouped into three hypotheses
$H\in\{\text{weak},\text{medium},\text{strong}\}$, and revises it with Bayes'
rule after every action I take.

The \textbf{prior} $\Pr(H)$ is the belief before I act; at the start of a hand it
is $(\Pr(\text{weak}),\Pr(\text{medium}),\Pr(\text{strong}))=(0.40,0.35,0.25)$.
The \textbf{likelihood} $\Pr(A\mid H)$ is how often each action $A$ (check, call,
bet, raise, fold) is taken with each strength; it encodes the intuition that a
raise is far likelier from a strong hand than a weak one. The default raise row,
for example, is $\Pr(\text{raise}\mid H)=(0.05,0.25,0.70)$. When I act, the bot
updates to the \textbf{posterior} by
\begin{equation}
\Pr(H\mid A)=\frac{\Pr(A\mid H)\,\Pr(H)}{\Pr(A)}
=\frac{\Pr(A\mid H)\,\Pr(H)}{\sum_{H'}\Pr(A\mid H')\,\Pr(H')},
\label{eq:bayes}
\end{equation}
where the denominator $\Pr(A)=\sum_{H'}\Pr(A\mid H')\Pr(H')$ is the
total-probability normalizer (the ``evidence'').

\example{Suppose I raise from the starting prior. Multiplying each likelihood by
its prior gives the unnormalized scores
\[
\text{weak}:0.05\times0.40=0.0200,\quad
\text{medium}:0.25\times0.35=0.0875,\quad
\text{strong}:0.70\times0.25=0.1750,
\]
which sum to $\Pr(\text{raise})=0.2825$. Normalizing yields the posterior
\[
\Pr(H\mid\text{raise})\approx(0.071,\ 0.310,\ 0.620).
\]
A single raise lifts the belief that I am strong from $25\%$ to about $62\%$ and
drops ``weak'' from $40\%$ to $7\%$.}

Each posterior becomes the prior for my next action, so a hand's belief is just a
running product of likelihoods that sharpens as I reveal more through my betting.
Crucially, this belief is the very distribution the Monte Carlo of Section~1 uses
to sample my hand---so when the bot reads me as strong, its estimated equity
drops automatically.

% ===========================================================================
\section{Learned Opponent Modeling: adapting across hands}
% ===========================================================================

The likelihood table above is a fixed guess about a \emph{generic} player. What
makes the bot interesting is that it \emph{learns} my personal likelihoods from
data. Every time a hand goes to showdown my cards are revealed, so the bot can
label my true strength tier and record which actions I took. Across many hands it
accumulates two counts per tier $T$: $n_T$, how many revealed hands were of that
tier, and $k_{a,T}$, how many of those involved action $a$.

The naive estimate $k_{a,T}/n_T$ is useless early on---a single hand would force a
probability to $0$ or $1$. So I smooth it with a \textbf{Beta prior}, treating
``did I take action $a$ on a tier-$T$ hand?'' as a Bernoulli with rate $\theta$
and putting a $\mathrm{Beta}(\alpha,\beta)$ prior on $\theta$ with $\alpha=2$ and
$\alpha+\beta=10$. Since the Beta is conjugate to the Bernoulli, the posterior
after the data is $\mathrm{Beta}(k_{a,T}+2,\,n_T-k_{a,T}+8)$, and its mean is the
likelihood the bot actually plugs into Equation~\eqref{eq:bayes}:
\begin{equation}
\Pr(a\mid T)=\frac{k_{a,T}+\alpha}{n_T+(\alpha+\beta)}
=\frac{k_{a,T}+2}{n_T+10}.
\label{eq:smooth}
\end{equation}
This formula is satisfying because it behaves correctly at both extremes. With no
data ($n_T=k_{a,T}=0$) every action starts at $2/10=0.20$, a neutral prior; as
$n_T\to\infty$ it converges to the empirical frequency $k_{a,T}/n_T$, so the prior
fades exactly as real evidence accumulates.

\example{Suppose I keep reaching showdown having raised with weak hands. After
$15$ such hands,
\[
\Pr(\text{raise}\mid\text{weak})=\frac{15+2}{15+10}=\frac{17}{25}=0.68,
\]
well above the $0.20$ default. Now a raise from me, run through
Equation~\eqref{eq:bayes}, gives $\Pr(\text{weak}\mid\text{raise})\approx0.64$
instead of the $\approx0.07$ the generic table produced. The bot has learned that
\emph{I} bluff-raise and reinterprets the identical action accordingly.}

So there are two timescales of inference working together: a fast belief update
\emph{within} a hand via Equation~\eqref{eq:bayes}, and a slow likelihood update
\emph{across} hands via Equation~\eqref{eq:smooth} that continually re-trains the
first.

% ===========================================================================
\section{Expected Value: choosing an action}
% ===========================================================================

Win probability alone can't decide a bet: a $30\%$ call is correct when the pot is
large, while an $80\%$ call can lose money at a bad price. The bot therefore ranks
actions by \textbf{expected value}, the average chip result $\mathbb{E}[O]=\sum_i
o_i\Pr(o_i)$, in a forward-looking form where money already in the pot is sunk and
only what I risk \emph{now} matters. With win probability $p$, loss probability
$q$, pot $P$, amount to call $c_{\text{call}}$, and an action of immediate cost
$\mathrm{cost}$ whose amount beyond a call is $\mathrm{extra}=\max(0,\,
\mathrm{cost}-c_{\text{call}})$, the unifying formula (ties are chip-neutral) is
\begin{equation}
\mathbb{E}[\text{action}]=p\,(P+\mathrm{extra})-q\cdot\mathrm{cost}.
\label{eq:ev}
\end{equation}
Specializing to each legal action,
\begin{align}
\mathbb{E}[\text{Fold}]&=0, &
\mathbb{E}[\text{Check}]&=p\,P, &
\mathbb{E}[\text{Call}]&=p\,P-q\,c_{\text{call}},\\
\mathbb{E}[\text{Bet}]&=p\,(P+\mathrm{cost})-q\cdot\mathrm{cost}, &
\mathbb{E}[\text{Raise}]&=p\,(P+\mathrm{extra})-q\cdot\mathrm{cost}. &&
\end{align}
Fold is the zero baseline, which is what makes folding ever correct: an action is
worth taking only if its EV clears $0$. The formula also reproduces standard
poker theory---requiring $\mathbb{E}[\text{Call}]=p\,P-q\,c_{\text{call}}>0$
rearranges to the familiar pot-odds rule
\begin{equation}
\frac{p}{q}>\frac{c_{\text{call}}}{P}.
\end{equation}
The bot then plays the maximizer, $a^\star=\argmax_a \mathbb{E}[a]$.

\example{Facing a \$10 bet into a \$30 pot, the simulation returns $p=0.35$,
$q=0.62$. A call has $\mathrm{cost}=c_{\text{call}}=10$, so
$\mathbb{E}[\text{Call}]=0.35(30)-0.62(10)=+\$4.3$; a raise to \$20 has
$\mathrm{cost}=20$ and $\mathrm{extra}=10$, so
$\mathbb{E}[\text{Raise}]=0.35(40)-0.62(20)=+\$1.6$; and
$\mathbb{E}[\text{Fold}]=0$. Since $\argmax\{0,\,4.3,\,1.6\}$ is the call, the bot
calls.}

This is where the whole pipeline closes: the $p$ and $q$ in Equation~\eqref{eq:ev}
come from the belief-weighted Monte Carlo, the belief comes from Bayesian updating,
and the likelihoods driving that update are themselves learned from past hands.

% ===========================================================================
\section*{Reflection}
% ===========================================================================

Building \emph{Probability Poker} changed how I think about probability. Coming in
as a first-year CS student, I knew the formulas in isolation, but implementing
them forced me to see how they depend on one another: a posterior is only as good
as its likelihoods, an equity estimate is only as good as the distribution it
samples from, and a decision is only as good as the probabilities feeding it.
Three lessons stuck with me. First, sampling is a genuinely powerful idea---the
$1/\sqrt N$ error bound meant I could trade a hopeless $10^9$-term sum for a few
thousand random trials and still act with confidence. Second, Bayesian updating
made ``learning'' concrete: priors are humility, likelihoods are evidence, and the
Beta-smoothed counts let the bot start cautious and grow certain only as data
earned it. Third, and most importantly, I learned that good decisions under
uncertainty are not about being sure---they are about combining imperfect beliefs
with the stakes of each choice, which is exactly what expected value does. The bot
never knows my cards, yet by chaining estimation, inference, learning, and
$\argmax$ it plays well anyway, and a little better every hand it watches.

\end{document}
